% Gerado em 26/11/2026 a partir do banco de produção (eventos bound_dispense e comandos dc1_motor).
% Só ensaios do ambiente Oficial (política "rigoroso"), 29/10 a 08/11/2026. Desde 14/11 o aparelho está no
% ambiente Simulado (banco tcc_simulado, política "apresentacao"); esses ensaios estão em ensaios_dispensacao_simulado.csv.
% Estilo igual às longtables da monografia (Capitulos/05_requisitos.tex).

{\small\setlength{\tabcolsep}{3pt}\begin{longtable}{p{0.035\textwidth}p{0.105\textwidth}p{0.04\textwidth}p{0.05\textwidth}p{0.155\textwidth}p{0.18\textwidth}p{0.13\textwidth}p{0.105\textwidth}}
\caption{Ensaios de dispensação integrada no submódulo SUB01 (aplicativo, API, central e ESP32)}\label{tab:ensaios-dispensacao}\\
\hline
\textbf{N.} & \textbf{Início} & \textbf{Qtd.} & \textbf{Ciclos} & \textbf{Sensor antes/depois} & \textbf{Resultado do aparelho} & \textbf{Observação humana} & \textbf{Estoque}\\\hline
\endfirsthead\hline
\textbf{N.} & \textbf{Início} & \textbf{Qtd.} & \textbf{Ciclos} & \textbf{Sensor antes/depois} & \textbf{Resultado do aparelho} & \textbf{Observação humana} & \textbf{Estoque}\\\hline
\endhead
1 & 29/10 13h34 & 1 & 0 & -- & Falha antes do ciclo (leitura do ADC alterada) & Saída de 1 (registrada em 07/11) & 13 $\rightarrow$ 12\\\hline
2 & 07/11 23h25 & 2 & 0 & -- & Recusa: estado inicial ambíguo & Sem saída & 12 $\rightarrow$ 12\\\hline
3 & 08/11 09h37 & 1 & 1 & 154 ocupado / 223 ocupado & Unidade retida & Sem saída & 12 $\rightarrow$ 12\\\hline
4 & 08/11 13h26 & 1 & 1 & 108,5 ocupado / 1916,5 vazio & Saída concluída & Saída de 1 & 12 $\rightarrow$ 11\\\hline
5 & 08/11 13h27 & 2 & 1 & 139 ocupado / 144 ocupado & Unidade retida (2 tentativas de desobstrução) & Não observado & 11 / pós ausente\\\hline
6 & 08/11 21h32 & 1 & 0 & -- & Recusa: estado inicial ambíguo & Sem saída & 11 $\rightarrow$ 11\\\hline
7 & 08/11 21h32 & 1 & 0 & -- & Recusa: estado inicial ambíguo & Sem saída & 11 $\rightarrow$ 11\\\hline
8 & 08/11 21h38 & 1 & 1 & 388 ocupado / -- & Ciclo interrompido; exige reconciliação & Saída de 1 & 11 $\rightarrow$ 10\\\hline
\end{longtable}}{\small\noindent Fonte: Autoria própria (2026).\par}


{\small\setlength{\tabcolsep}{3pt}\begin{longtable}{p{0.65\textwidth}p{0.25\textwidth}}
\caption{Qualidade e elegibilidade dos oito registros do Oficial}\label{tab:indicadores-dispensacao}\\
\hline\textbf{Critério de leitura} & \textbf{Registros}\\\hline\endfirsthead
\hline\textbf{Critério de leitura} & \textbf{Registros}\\\hline\endhead
Total de linhas no CSV & 8\\\hline
Quantidade observada preenchida & 7/8\\\hline
Observação tardia incompatível com o ciclo & 1 (ensaio 1)\\\hline
Sem quantidade observada e sem estoque posterior & 1 (ensaio 5)\\\hline
Linhas com quantidade e estoque antes/depois, excluído o caso tardio & 6 (2, 3, 4, 6, 7 e 8)\\\hline
Nessas seis, redução do estoque registrado igual à quantidade relatada & 6/6; consistência contábil\\\hline
Saídas unitárias relatadas sem a inconsistência temporal do ensaio 1 & 2 (4 e 8)\\\hline
Recusas antes de mover com relato de zero saída & 3 (2, 6 e 7)\\\hline
Ciclo retido com relato de zero saída & 1 (3)\\\hline
\end{longtable}\noindent Fonte: Autoria própria (2026).\par}
